\chapter{Forwarding toyRISC}

\section{Structure}

{\small
\begin{shaded*}
\begin{lstlisting}[language=Verilog, caption= , morekeywords={always_comb, always_ff, logic, loop, repeat, doInParallel, for}]
/************************************************************************
File name: toyRISCsystem.sv
Circuit name:
Description:
************************************************************************/
`include "0_DEFINES.vh"
module toyRISCsystem(   output  logic   inta    ,
                        input   logic   intIn   ,
                                        reset   ,
                                        clk     );
    logic   [31:0]  instruction ;
    logic   [9:0]   PC          ;
    logic   [31:0]  dataIn      ;
    logic   [31:0]  dataOut     ;
    logic   [9:0]   dataAddr    ;
    logic           dataRead    ;
    logic           dataWrite   ;

    toyRISC processor(  instruction ,
                        PC          ,
                        intIn       ,
                        inta        ,
                        dataIn      ,
                        dataOut     ,
                        dataAddr    ,
                        dataRead    ,
                        dataWrite   ,
                        reset       ,
                        clk         );
    memorySystem memSys(instruction ,
                        PC          ,
                        dataIn      ,
                        dataOut     ,
                        dataAddr    ,
                        dataRead    ,
                        dataWrite   ,
                        clk         );
endmodule
\end{lstlisting}
\end{shaded*}
}

{\small
\begin{shaded*}
\begin{lstlisting}[language=Verilog, caption= , morekeywords={always_comb, always_ff, logic, loop, repeat, doInParallel, for}]
/************************************************************************
File name: memorySystem.sv
Circuit name:
Description:
************************************************************************/
module memorySystem(output  logic   [31:0]  instruction ,
                    input   logic   [9:0]   PC          ,
                    output  logic   [31:0]  dataIn      ,
                    input   logic   [31:0]  dataOut     ,
                    input   logic   [9:0]   dataAddr    ,
                    input   logic           dataRead    ,
                    input   logic           dataWrite   ,
                    input   logic           clk         );

    logic   [31:0]  dataMemory[0:1023]  ;
    logic   [31:0]  progMemory[0:1023]  ;

    assign instruction = progMemory[PC] ;

    always_ff @(posedge clk) begin
        if (dataWrite) dataMemory[dataAddr] <= dataOut  ;
    end

    assign dataIn   = dataMemory[dataAddr];

endmodule
\end{lstlisting}
\end{shaded*}
}

{\small
\begin{shaded*}
\begin{lstlisting}[language=Verilog, caption= , morekeywords={always_comb, always_ff, logic, loop, repeat, doInParallel, for}]
/************************************************************************
File name: .sv
Circuit name:
Description:
************************************************************************/
/*************************************************************************
File name: toyRISC.sv
                        toyRISC
*************************************************************************/
`include "0_DEFINES.vh"
module toyRISC( input   logic [31:0]  instruction   ,
                output  logic [9:0]   pc            ,
                input   logic         intIn         ,
                output  logic         inta          ,
                input   logic [31:0]  dataIn        ,
                output  logic [31:0]  dataOut       ,
                output  logic [9:0]   dataAddr      ,
                output  logic         dataRead      ,
                output  logic         dataWrite     ,
                input   logic         reset         ,
                input   logic         clk           );
    logic [31:0]    instruction1;
    logic [31:0]    instruction2;
    logic [31:0]    instruction3;
    logic [31:0]    instruction4;
    logic [9:0]     pc1         ;
    logic [31:0]    leftOp      ;
    logic [1:0]     nextPCsel   ;
    logic           we4         ;
    logic [31:0]    result4     ;
    logic [3:0]     opSel       ;
    logic [3:0]     opSel2      ;
    logic           zero        ;
    logic [31:0]    leftOp2     ;
    logic [31:0]    rightOp2    ;
    logic [31:0]    leftOp3     ;
    logic [31:0]    rightOp3    ;
    logic [31:0]    result3     ;
    logic           we          ;
    logic           memSel      ;

    DECODE dcd( .instruction1   (instruction1   ),
                .instruction2   (instruction2   ),
                .instruction3   (instruction3   ),
                .zero           (zero           ),
                .intIn          (intIn          ),

                .inta           (inta           ),
                .nextPCsel      (nextPCsel      ),
                .opSel          (opSel          ),
                .memSel         (memSel         ),
                .we             (we             ),
                .dataRead       (dataRead       ),
                .dataWrite      (dataWrite      ),

                .reset          (reset          ),
                .clk            (clk            ));

    INSTRFETCH progCount(   instruction ,
                            leftOp      ,
                            nextPCsel   ,

                            pc          ,
                            instruction1,
                            pc1         ,

                            reset       ,
                            clk         );

    OPSFETCH regFile(   instruction1,
                        instruction4,
                        pc1         ,
                        inta        ,
                        we4         ,
                        result4     ,
                        opSel       ,

                        zero        ,
                        leftOp      ,
                        leftOp2     ,
                        rightOp2    ,
                        opSel2      ,
                        instruction2,

                        clk         );

    EXECUTE execute(leftOp2     ,
                    rightOp2    ,
                    opSel2      ,
                    instruction2,
                    result4     ,

                    leftOp3     ,
                    result3     ,
                    rightOp3    ,
                    instruction3,

                    clk         );

    DATAMEM datamem(result3     ,
                    instruction3,
                    we          ,
                    memSel      ,
                    dataIn      ,

                    result4     ,
                    instruction4,
                    we4         ,

                    clk         );

    assign dataAddr = leftOp3   ;
    assign dataOut  = rightOp3  ;
endmodule // Synthesis results: #LUT=335, #FF=204, #DSP=3
\end{lstlisting}
\end{shaded*}
}

{\small
\begin{shaded*}
\begin{lstlisting}[language=Verilog, caption= , morekeywords={always_comb, always_ff, logic, loop, repeat, doInParallel, for}]
/************************************************************************
File name: DECODE.sv
Circuit name:
Description:
************************************************************************/
`include "0_DEFINES.vh"
module DECODE(  input   logic [31:0]    instruction1,
                input   logic [31:0]    instruction2,
                input   logic [31:0]    instruction3,
                input   logic           zero        ,
                input   logic           intIn       ,

                output  logic           inta        ,
                output  logic [1:0]     nextPCsel   ,
                output  logic [3:0]     opSel       ,
                output  logic           memSel      ,
                output  logic           we          ,
                output  logic           dataRead    ,
                output  logic           dataWrite   ,

                input   logic           reset       ,
                input   logic           clk         );

    logic [5:0] opCode1     ;
    logic [4:0] destAddr1   ;
    logic [4:0] leftAddr1   ;
    logic [4:0] rightAddr1  ;
    logic [4:0] destAddr2   ;
    logic [5:0] opCode3     ;
    logic [4:0] destAddr3   ;
    logic       leftFwd     ;
    logic       rightFwd    ;

    assign opCode1      = instruction1[31:26]   ;
    assign leftAddr1    = instruction1[20:16]   ;
    assign rightAddr1   = instruction1[15:11]   ;
    assign destAddr2    = instruction2[25:21]   ;
    assign destAddr3    = instruction3[25:21]   ;
    assign opCode3      = instruction3[31:26]   ;


// Interrupt automaton
    logic [2:0] intState    ;
/*
    intState = 000: intDisable
    intState = 001: intEnable
    intState = 010: intExec1
    intState = 011: intExec2
    intState = 100: intExec3
    intState = 101: intExec4
*/
    always_ff @(posedge clk)
        if (reset)                              intState <= 3'b000  ;
         else
            begin
            if (opCode1 == `eint)               intState <= 3'b001  ;
            if (opCode1 == `dint)               intState <= 3'b000  ;
            if (intIn && (intState == 3'b001))  intState <= 3'b010  ;
            if (intState == 3'b010)             intState <= 3'b011  ;
            if (intState == 3'b011)             intState <= 3'b100  ;
            if (intState == 3'b100)             intState <= 3'b101  ;
            if (intState == 3'b101)             intState <= 3'b000  ;
            end
    assign inta = intState == 3'b011    ;
// end Interrupt automaton

// Program control
    logic [5:0] jmpSel  ;

    assign jmpSel = (intState > 3'b011) ? `nop : opCode1    ;

    always_comb if (inta)         nextPCsel = 2'b11                 ;
                 else
                    case(jmpSel)
                        `halt   : nextPCsel = 2'b00                 ;
                        `rjmp   : nextPCsel = 2'b10                 ;
                        `zbr    : nextPCsel = zero ? 2'b10 : 2'b01  ;
                        `nzbr   : nextPCsel = zero ? 2'b01 : 2'b10  ;
                        `ret    : nextPCsel = 2'b11                 ;
                        default : nextPCsel = 2'b01                 ;
                    endcase
// end prorgram control

// Forwarding section
    assign leftFwdY  =
        (leftAddr1  == destAddr2) && (opCode1[5] == 1'b1);
    assign rightFwdY =
        (rightAddr1 == destAddr2) && (opCode1[5] == 1'b1);
    assign leftFwdO  =
        (leftAddr1  == destAddr3) && (opCode1[5] == 1'b1);
    assign rightFwdO =
        (rightAddr1 == destAddr3) && (opCode1[5] == 1'b1);
    always_comb
        case(opCode1)
            `addv   :   opSel[1:0] = 2'b11  ;
            `multv  :   opSel[1:0] = 2'b11  ;
            `addvc  :   opSel[1:0] = 2'b11  ;
            `val    :   opSel[1:0] = 2'b11  ;
            default :   if (rightFwdY)          opSel[1:0] = 2'b01  ;
                         else if (rightFwdO)    opSel[1:0] = 2'b10  ;
                               else             opSel[1:0] = 2'b00  ;
        endcase
    always_comb
        if (leftFwdY)       opSel[3:2] = 2'b01  ;
         else if (leftFwdO) opSel[3:2] = 2'b10  ;
               else         opSel[3:2] = 2'b00  ;
// end forwarding section

    assign dataWrite    = opCode3 == `store                         ;
    assign dataRead     = opCode3 == `read                          ;
    assign memSel       = opCode3 == `read                          ;
    assign we           = ((opCode3[5:4] == 2'b11)  |
                           (opCode3 == `val)        |
                           (opCode3 == `read)) & !(intState > 3'b001);
endmodule
\end{lstlisting}
\end{shaded*}
}

{\small
\begin{shaded*}
\begin{lstlisting}[language=Verilog, caption= , morekeywords={always_comb, always_ff, logic, loop, repeat, doInParallel, for}]
/************************************************************************
File name: INSTRFETCH.sv
Circuit name:
Description:
************************************************************************/
module INSTRFETCH ( input   logic [31:0]    instruction ,
                    input   logic [31:0]    leftOp      ,
                    input   logic [1:0]     nextPCsel   ,

                    output  logic [9:0]     pc          ,
                    output  logic [31:0]    instruction1,
                    output  logic [9:0]     pc1         ,

                    input   logic           reset       ,
                    input   logic           clk         );

    // PipeReg0
    always_ff @(posedge clk)
        if (reset)  begin       pc <= -1                        ;
                                instruction1 <= 0               ;
                    end
         else begin case(nextPCsel)
                        2'b00:  pc <= pc                        ;
                        2'b01:  pc <= pc + 1                    ;
                        2'b10:  pc <= instruction1[9:0] + pc1   ;
                        2'b11:  pc <= leftOp                    ;
                    endcase
    // PipeReg1
                        instruction1 <= instruction ;
                        pc1          <= pc          ;
                end
endmodule
\end{lstlisting}
\end{shaded*}
}

{\small
\begin{shaded*}
\begin{lstlisting}[language=Verilog, caption= , morekeywords={always_comb, always_ff, logic, loop, repeat, doInParallel, for}]
/************************************************************************
File name: OPSFETCH.sv
Circuit name:
Description:
************************************************************************/
module OPSFETCH(    input   logic [31:0]    instruction1,
                    input   logic [31:0]    instruction4,
                    input   logic [9:0]     pc1         ,
                    input   logic           inta        ,
                    input   logic           we4         ,
                    input   logic [31:0]    result4     ,
                    input   logic [3:0]     opSel       ,

                    output  logic           zero        ,
                    output  logic [31:0]    leftOp      ,
                    output  logic [31:0]    leftOp2     ,
                    output  logic [31:0]    rightOp2    ,
                    output  logic [3:0]     opSel2      ,
                    output  logic [31:0]    instruction2,

                    input   logic           clk         );

    logic   [31:0]  rf[0:31]    ;
    logic   [31:0]  rightOp     ;
    logic           we          ;
    logic   [4:0]   destAddr    ;
    logic   [31:0]  result      ;
    logic   [4:0]   leftAddr    ;
    logic   [4:0]   rightAddr   ;


    always_ff @(posedge clk) if (we) rf[destAddr] <= result ;

    assign leftOp       =
        ((destAddr == leftAddr) && we) ? result : rf[leftAddr]    ;
    assign rightOp      =
        ((destAddr == rightAddr) && we) ? result : rf[rightAddr]  ;
    assign zero         = leftOp == 0           ;
    assign rightAddr    = instruction1[15:11]   ;

    // PipeReg2:
    always_ff @(posedge clk)
        begin
            leftOp2         <=  leftOp          ;
            rightOp2        <=  rightOp         ;
            opSel2          <=  opSel           ;
            instruction2    <=  instruction1    ;
        end

    intUnit intunit(.we         (we                     ),
                    .destAddr   (destAddr               ),
                    .result     (result                 ),
                    .leftAddr   (leftAddr               ),

                    .we4        (we4                    ),
                    .destAddr4  (instruction4[25:21]    ),
                    .result4    (result4                ),
                    .leftAddr1  (instruction1[20:16]    ),
                    .pc1        (pc1                    ),
                    .inta       (inta                   ));
endmodule
\end{lstlisting}
\end{shaded*}
}

{\small
\begin{shaded*}
\begin{lstlisting}[language=Verilog, caption= , morekeywords={always_comb, always_ff, logic, loop, repeat, doInParallel, for}]
/************************************************************************
File name: EXECUTE.sv
Circuit name:
Description:
************************************************************************/
`include "0_DEFINES.vh"
module EXECUTE( input   logic [31:0]    leftOp2     ,
                input   logic [31:0]    rightOp2    ,
                input   logic [3:0]     opSel2      ,
                input   logic [31:0]    instruction2,
                input   logic [31:0]    result4     ,

                output  logic [31:0]    leftOp3     ,
                output  logic [31:0]    result3     ,
                output  logic [31:0]    rightOp3    ,
                output  logic [31:0]    instruction3,

                input                   clk         );

    logic [31:0]    rightIn ;
    logic [31:0]    leftIn  ;
    logic [31:0]    out     ;
    logic [31:0]    value   ;

    assign value = {{16{instruction2[15]}}, instruction2[15:0]} ;

    always_comb case(opSel2[1:0])
                    2'b00: rightIn = rightOp2   ;
                    2'b01: rightIn = result3    ;
                    2'b10: rightIn = result4    ;
                    2'b11: rightIn = value      ;
                endcase
    always_comb case(opSel2[3:2])
                    2'b00: leftIn = leftOp2;
                    2'b01: leftIn = result3;
                    2'b10: leftIn = result4;
                    2'b11: leftIn = 31'b0   ;
                endcase

    ALU alu(.out    (out                ),
            .leftIn (leftIn             ),
            .rightIn(rightIn            ),
            .opCode (instruction2[31:26]));

    // PipeReg3
    always_ff @(posedge clk)
        begin   leftOp3         <= leftIn       ;
                result3         <= out          ;
                rightOp3        <= rightIn      ;
                instruction3    <= instruction2 ;
        end
endmodule
\end{lstlisting}
\end{shaded*}
}

{\small
\begin{shaded*}
\begin{lstlisting}[language=Verilog, caption= , morekeywords={always_comb, always_ff, logic, loop, repeat, doInParallel, for}]
/************************************************************************
File name: DATAMEM.sv
Circuit name:
Description:
************************************************************************/
module DATAMEM( input   logic [31:0]    result3     ,
                input   logic [31:0]    instruction3,
                input   logic           we          ,
                input   logic           memSel      ,
                input   logic [31:0]    dataIn      ,

                output  logic [31:0]    result4     ,
                output  logic [31:0]    instruction4,
                output  logic           we4         ,

                input   logic           clk         );

    always_ff @(posedge clk)
        begin   result4         <= memSel ? dataIn : result3;
                instruction4    <= instruction3             ;
                we4             <= we                       ;
        end
endmodule
\end{lstlisting}
\end{shaded*}
}


{\small
\begin{shaded*}
\begin{lstlisting}[language=Verilog, caption= , morekeywords={always_comb, always_ff, logic, loop, repeat, doInParallel, for}]
/************************************************************************
File name: ALU.sv
Circuit name:
Description:
************************************************************************/
`include "0_DEFINES.vh"
module ALU( output  logic   [31:0]  out     ,
            input   logic   [31:0]  leftIn  ,
                                    rightIn ,
            input   logic   [5:0]   opCode  );

    logic [32:0]    sum ;
    logic [32:0]    dif ;

    assign sum = leftIn + rightIn   ;
    assign dif = leftIn - rightIn   ;

    always_comb case(opCode)
                    `add    : out = sum[31:0]                       ;
                    `sub    : out = dif[31:0]                       ;
                    `addv   : out = sum[31:0]                       ;
                    `mult   : out = leftIn * rightIn                ;
                    `multv  : out = leftIn * rightIn                ;
                    `addc   : out = sum[32]                         ;
                    `subc   : out = dif[32]                         ;
                    `addvc  : out = sum[32]                         ;
                    `lsh    : out = leftIn >> 1                     ;
                    `ash    : out = {leftIn[31], leftIn[31:1]}      ;
                    `move   : out = leftIn                          ;
                    `swap   : out = {leftIn[15:0], leftIn[31:16]}   ;
                    `bwnot  : out = ~leftIn                         ;
                    `bwand  : out = leftIn & rightIn                ;
                    `bwor   : out = leftIn | rightIn                ;
                    `bwxor  : out = leftIn ^ rightIn                ;
                    //`load : out = rightIn                         ;
                    `val    : out = rightIn                         ;
                    default : out = leftIn                          ;
                endcase
endmodule
\end{lstlisting}
\end{shaded*}
}

{\small
\begin{shaded*}
\begin{lstlisting}[language=Verilog, caption= , morekeywords={always_comb, always_ff, logic, loop, repeat, doInParallel, for}]
/************************************************************************
File name: intUnit.sv
Circuit name:
Description:
************************************************************************/
module intUnit( output  logic           we          ,
                output  logic   [4:0]   destAddr    ,
                output  logic   [31:0]  result      ,
                output  logic   [4:0]   leftAddr    ,

                input   logic           we4         ,
                input   logic   [4:0]   destAddr4   ,
                input   logic   [31:0]  result4     ,
                input   logic   [4:0]   leftAddr1   ,
                input   logic   [9:0]   pc1         ,
                input   logic           inta        );

    assign we       = inta ? 1'b1           : we4       ;
    assign destAddr = inta ? 5'b11110       : destAddr4 ;
    assign result   = inta ? {21'b0, pc1}   : result4   ;
    assign leftAddr = inta ? 5'b11111       : leftAddr1 ;
endmodule
\end{lstlisting}
\end{shaded*}
}

\section{Code generator}


{\small
\begin{shaded*}
\begin{lstlisting}[language=Verilog, caption= , morekeywords={always_comb, always_ff, logic, loop, repeat, doInParallel, for}]
/************************************************************************
File name: RISCcodeGenerator.sv
Circuit name:
Description:
************************************************************************/
    reg [5:0]   opCode          ;
    reg [4:0]   d               ;
    reg [4:0]   l               ;
    reg [4:0]   r               ;
    reg [15:0]  v               ;
    reg [9:0]   addrCounter     ;
    reg [9:0]   labelTab[0:1023];

    `include "0_DEFINES.vh"

    task endLine;
     begin
        dut.memSys.progMemory[addrCounter][31:0] =
            {   opCode  ,
                d       ,
                l       ,
                v       }                   ;
            addrCounter = addrCounter + 1   ;
     end
    endtask

    // sets labelTab in the first pass
    // associating 'counter' with 'labelIndex'
    task LB ;
        input [5:0] labelIndex;

        labelTab[labelIndex] = addrCounter;
    endtask
    // uses the content of labelTab in the second pass
    task ULB;
        input [5:0] labelIndex;

        v = labelTab[labelIndex] - addrCounter  ;
        endtask

// CONTROL INSTRUCTIONS
    task NOP; // no operation
        begin   opCode  = `nop  ;
                {d,l,v} = 26'b0 ;
                endLine         ;
        end
    endtask

    task RJMP; // relative jump
        input   [15:0]  label   ;

        begin   opCode  = `rjmp ;
                d       = 5'b0  ;
                l       = 5'b0  ;
                ULB(label)      ;
                endLine         ;
        end
    endtask

    task BRZ; // branch if zero
        input   [4:0]   left    ;
        input   [15:0]  label   ;

        begin   opCode  = `zbr  ;
                d       = 5'b0  ;
                l       = left  ;
                ULB(label)      ;
                endLine         ;
        end
    endtask

    task BRNZ; // branch if not zero
        input   [4:0]   left    ;
        input   [15:0]  label   ;

        begin   opCode  = `nzbr ;
                d       = 5'b0  ;
                l       = left  ;
                ULB(label)      ;
                endLine         ;
        end
    endtask

    task RET; // return from subroutine
        input   [4:0] left  ;

        begin   opCode  = `ret  ;
                d       = 5'b0  ;
                l       = left  ;
                v       = 16'b0 ;
                endLine         ;
        end
    endtask

    task HALT; // halt running
        begin   opCode  = `halt ;
                {d,l,v} = 26'b0 ;
                endLine         ;
        end
    endtask

    task EI; // enable interrupt
        begin   opCode  = `eint ;
                {d,l,v} = 26'b0 ;
                endLine         ;
        end
    endtask

    task DI; // disable interrupt
        begin   opCode  = `dint ;
                {d,l,v} = 26'b0 ;
                endLine         ;
        end
    endtask

// ARITHMETIC & LOGIC INSTRUCTIONS
    task ADD; // addition
        input   [4:0]   dest    ;
        input   [4:0]   left    ;
        input   [4:0]   right   ;

        begin   opCode  = `add          ;
                d       = dest          ;
                l       = left          ;
                v       = {right, 11'b0};
                endLine                 ;
        end
    endtask

    task SUB; // subtract
        input   [4:0]   dest    ;
        input   [4:0]   left    ;
        input   [4:0]   right   ;

        begin   opCode  = `sub          ;
                d       = dest          ;
                l       = left          ;
                v       = {right, 11'b0};
                endLine                 ;
        end
    endtask

    task ADDV; // addition with value
        input   [4:0]   dest    ;
        input   [4:0]   left    ;
        input   [15:0]  value   ;

        begin   opCode  = `addv ;
                d       = dest  ;
                l       = left  ;
                v       = value ;
                endLine         ;
        end
    endtask

    task MULT; // multiplication
        input   [4:0]   dest    ;
        input   [4:0]   left    ;
        input   [4:0]   right   ;

        begin   opCode  = `mult         ;
                d       = dest          ;
                l       = left          ;
                v       = {right, 11'b0};
                endLine                 ;
        end
    endtask

    task MULTV; // multiplication with value
        input   [4:0]   dest    ;
        input   [4:0]   left    ;
        input   [15:0]  value   ;

        begin   opCode  = `multv;
                d       = dest  ;
                l       = left  ;
                v       = value ;
                endLine         ;
        end
    endtask

    task ADDC; // carry from addition
        input   [4:0]   dest    ;
        input   [4:0]   left    ;
        input   [4:0]   right   ;

        begin   opCode  = `addc         ;
                d       = dest          ;
                l       = left          ;
                v       = {right, 11'b0};
                endLine                 ;
        end
    endtask

    task SUBC; // carry from subtract
        input   [4:0]   dest    ;
        input   [4:0]   left    ;
        input   [4:0]   right   ;

        begin   opCode  = `subc         ;
                d       = dest          ;
                l       = left          ;
                v       = {right, 11'b0};
                endLine                 ;
        end
    endtask

    task ADDVC; // carry from addition with value
        input   [4:0]   dest    ;
        input   [4:0]   left    ;
        input   [15:0]  value   ;

        begin   opCode  = `addvc;
                d       = dest  ;
                l       = left  ;
                v       = value ;
                endLine         ;
        end
    endtask

    task LSH; // logic shift with one position
        input   [4:0]   dest    ;
        input   [4:0]   left    ;

        begin   opCode  = `lsh  ;
                d       = dest  ;
                l       = left  ;
                v       = 16'b0 ;
                endLine         ;
        end
    endtask

    task ASH; // arithmetic shift with one porition
        input   [4:0]   dest    ;
        input   [4:0]   left    ;

        begin   opCode  = `ash  ;
                d       = dest  ;
                l       = left  ;
                v       = 16'b0 ;
                endLine         ;
        end
    endtask

    task MOVE; // data move inside register file
        input   [4:0]   dest    ;
        input   [4:0]   left    ;

        begin   opCode  = `move ;
                d       = dest  ;
                l       = left  ;
                v       = 16'b0 ;
                endLine         ;
        end
    endtask

    task SWAP; // swap in register
        input   [4:0]   dest    ;
        input   [4:0]   left    ;

        begin   opCode  = `swap ;
                d       = dest  ;
                l       = left  ;
                v       = 16'b0 ;
                endLine         ;
        end
    endtask

    task NOT; // bitwise NOT
        input   [4:0]   dest    ;
        input   [4:0]   left    ;

        begin   opCode  = `bwnot;
                d       = dest  ;
                l       = left  ;
                v       = 16'b0 ;
                endLine         ;
        end
    endtask

    task AND; // bitwise AND
        input   [4:0]   dest    ;
        input   [4:0]   left    ;
        input   [4:0]   right   ;

        begin   opCode  = `bwand        ;
                d       = dest          ;
                l       = left          ;
                v       = {right, 11'b0};
                endLine                 ;
        end
    endtask

    task OR; // bitwise OR
        input   [4:0]   dest    ;
        input   [4:0]   left    ;
        input   [4:0]   right   ;

        begin   opCode  = `bwor     ;
                d       = dest          ;
                l       = left          ;
                v       = {right, 11'b0};
                endLine                 ;
        end
    endtask

    task XOR; // bitwise XOR
        input   [4:0]   dest    ;
        input   [4:0]   left    ;
        input   [4:0]   right   ;

        begin   opCode  = `bwxor        ;
                d       = dest          ;
                l       = left          ;
                v       = {right, 11'b0};
                endLine                 ;
        end
    endtask

// DATA TRANSFER INSTRUCTIONS
    task READ; // data read
        input   [4:0]   dest    ;
        input   [4:0]   left    ;

        begin   opCode  = `read ;
                d       = dest  ;
                l       = left  ;
                v       = 16'b0 ;
                endLine         ;
        end
    endtask

    task STORE; // data store
        input   [4:0]   left    ;
        input   [4:0]   right   ;

        begin   opCode  = `store        ;
                d       = 5'b0          ;
                l       = left          ;
                v       = {right, 11'b0};
                endLine                 ;
        end
    endtask

    task VAL; // value load
        input   [4:0]   dest    ;
        input   [15:0]  value   ;

        begin   opCode  = `val  ;
                d       = dest  ;
                l       = 5'b0  ;
                v       = value ;
                endLine         ;
        end
    endtask

    // RUNNING
    initial begin   addrCounter = 0;
                    `include "0_program.sv" // first pass
                    addrCounter = 0;
                    `include "0_program.sv" // second pass
            end
\end{lstlisting}
\end{shaded*}
}


\section{Simulator}

{\small
\begin{shaded*}
\begin{lstlisting}[language=Verilog, caption= , morekeywords={always_comb, always_ff, logic, loop, repeat, doInParallel, for}]
/************************************************************************
File name: testRISC.sv
Circuit name:
Description:
************************************************************************/
`include "0_DEFINES.vh"
module testRISC;
    logic   inta    ;
    logic   intIn   ;
    logic   reset   ;
    logic   clk     ;

    integer i   ;

    `include "0_RISCcodeGenerator.sv"

    initial begin               clk = 0     ;
                    forever #1  clk = ~clk  ;
            end

    initial
        begin           intIn       = 1 ;
                        reset       = 1 ;
                        // DISPLAY THE CONTENT OF PROGRAM MEMORY
                        for (i=0; i<16; i=i+1)
                            $display("progMemory[%0d] \t = %b", i,
                                    dut.memSys.progMemory[i])   ;
                #4      reset       = 0 ;
                #100    $finish     ;
        end

    toyRISCsystem dut(  inta    ,
                        intIn   ,
                        reset   ,
                        clk     );

    // MONITOR FOR PROGRAM LAOD & CONTROLLER
    initial begin
     $monitor("t=%0d pc=%d  RF=[%0d, %0d, %0d, %0d, %0d, %0d, %0d, %0d]
              dataWrite=%0b DM=[%0d, %0d, %0d, %0d] intState=%b inta=%b",
                $time,
                dut.processor.pc,
                dut.processor.regFile.rf[0],
                dut.processor.regFile.rf[1],
                dut.processor.regFile.rf[2],
                dut.processor.regFile.rf[3],
                dut.processor.regFile.rf[4],
                dut.processor.regFile.rf[5],
                dut.processor.regFile.rf[30],
                dut.processor.regFile.rf[31],
                dut.processor.dataWrite,
                dut.memSys.dataMemory[0],
                dut.memSys.dataMemory[1],
                dut.memSys.dataMemory[2],
                dut.memSys.dataMemory[3],
                dut.processor.dcd.intState,
                dut.processor.dcd.inta);
    end
endmodule
\end{lstlisting}
\end{shaded*}
}



\section{Testing}

{\small
\begin{shaded*}
\begin{lstlisting}[language=Verilog, caption= , morekeywords={always_comb, always_ff, logic, loop, repeat, doInParallel, for}]
/************************************************************************
File name: program.sv
Circuit name:
Description:
************************************************************************/
/*
            NOP         ;
            VAL(0,1)    ;
            VAL(1,2)    ;
            VAL(2,3)    ;
            VAL(3,4)    ;
            VAL(4,5)    ;
            //NOP           ;
            //NOP           ;
            //NOP           ;
            ADD(5,4,3)  ;
            HALT        ;
            HALT        ;
//*/


/*
            VAL(0,2)    ;
            VAL(1,4)    ;
            VAL(2,8)    ;
            MULT(3,1,2) ;
            ADDV(0,0,5) ;
            NOP         ;
            HALT        ;
            HALT        ;
//*/
/*
            VAL(31,13)  ;
            VAL(2,23)   ;
            VAL(0,24)   ;
            VAL(1,13)   ;
            EI          ;
            ADDV(2,31,1);
            VAL(1,1)    ;
            VAL(2,222)  ;
            NOP         ;
            ADD(0,1,2)  ;
            HALT        ;
            HALT        ;
            NOP         ;
        // subroutine triggered by interrupt
            ADDV(30,30,-2)  ;
            VAL(3,44)       ;
            NOP             ;
            NOP             ;
            NOP             ;
            RET(30)         ;
            NOP             ;
//*/
//*
            NOP         ;
            VAL(0,5)    ;
    LB(1);  ADDV(0,0,-1);
            NOP         ;
            NOP         ;
            BRNZ(0,1)   ;
            NOP         ;
            HALT        ;
            HALT        ;
//*/
/*
            VAL(1,55)   ;
            VAL(0,1)    ;
            STORE(0,1)  ;
            NOP         ;
            NOP         ;
            READ(2,0)   ;
            HALT        ;
            HALT        ;
//*/











\end{lstlisting}
\end{shaded*}
}

